First, we note that the reduced costs corresponding to the basic
variables are zero. Therefore, $\bar{c}_1 = \bar{c}_2=0$.

The reduced cost for the non basic variable $x_3$ is defined as
\[
\bar{c}_3 = c_3 - c_B^TB^{-1}A_3 = c_3 + c_B^T (d_3)_B,
\]
where $c_3$ is the coefficient of $x_3$ in the objective function
($c_3=0$), $c_B$ are the entries of the cost vector corresponding
to the basic variables:
\[
c_B = \cvect{1 \\ 0},
\]
and $(d_3)_B$ are the basic entries of the basic direction $d_3$:
\[
(d_3)_B = -B^{-1} A_3 = -\begin{pmatrix*}[r]
\frac{2}{3}&-\frac{1}{3}\\
\frac{1}{3} & \frac{1}{3}
\end{pmatrix*} \cvect{1 \\ 0} = \cvect{-\frac{2}{3} \\ -\frac{1}{3}}.
\]
Therefore,
\[
\bar{c}_3 = c_B^T (d_3)_B = -\frac{2}{3}.
\]
Note that it is negative. It means that the objective function is
decreasing along that direction.  Each step of size 1 into this
direction will decrease the value of the objective function by $2/3$.

The reduced cost for the non basic variable $x_4$ is defined as
\[
\bar{c}_4 = c_4 - c_B^TB^{-1}A_4 = c_4 + c_B^T (d_4)_B,
\]
where $c_4$ is the coefficient of $x_4$ in the objective function
($c_4=0$), $c_B$ are the entries of the cost vector corresponding
to the basic variables:
\[
c_B = \cvect{1 \\ 0},
\]
and $(d_4)_B$ are the basic entries of the basic direction $d_4$:
\[
(d_4)_B = -B^{-1} A_4 = -\begin{pmatrix*}[r]
\frac{2}{3}&-\frac{1}{3}\\
\frac{1}{3} & \frac{1}{3}
\end{pmatrix*} \cvect{0 \\ 1} = \cvect{\frac{1}{3} \\ -\frac{1}{3}}.
\]
Therefore,
\[
\bar{c}_4 = c_B^T (d_4)_B = \frac{1}{3}.
\]
Note that it is positive. It means that the objective function is
increasing along that direction.   Each step of size 1 into this
direction will increase the value of the objective function by $1/3$.

As a final note, it happens that the point $(0,1,0,0)$ is an optimal
solution of the problem. Indeed, it is feasible, and $x_1$ is at its
minimum value. Still, not all reduced costs are non negative. But the
direction $d_3$ corresponding to the negative reduced cost is
infeasible (as seen in the previous exercise). It
is an example that shows that the condition ``all reduced costs non
negative'' is a sufficient condition for optimality, but not a
necessary condition. In particular, Theorem 6.30 does not apply as the
feasible basic solution is degenerate.
